\begin{abstract}
Multi-teacher on-policy distillation (MOPD) aims to combine capabilities
from several teachers in one student. Successful transfer from each teacher
separately does not establish that their combined updates improve every
task. We analyze joint progress by separating the expected effect of an
update on each task from uncertainty due to finite sampling. For losses on
fixed student-generated states and a frozen optimizer scaling, we bound the
probability of any non-descent first-order direction when the mean update
improves every task to first order. The analysis distinguishes harmful mean
interactions, which allocation at fixed weights cannot correct in this local
model, from sampling errors in otherwise useful directions.
Gradient-Preconditioned Adaptive Sampling (GPAS) minimizes an estimated
total-variance criterion; an
analytical example shows why this need not minimize uncertainty in task
progress. We specify a four-domain study with separate-distillation
references, controlled updates from common checkpoints, and repeated
training comparisons. The evaluation distinguishes fixed-state loss,
fresh-policy loss, downstream capability, and computation.
\missing{Add the observed integration gap, diagnostic prediction accuracy,
and measured capability and compute effects after completing the study.}
\end{abstract}
